\chapter{DE LA GRAMMAIRE À LA LINGUISTIQUE}
\label{chapter-3}
\markright{De la grammaire à la linguistique}
\origpage{58}{044}
\begin{keybox}[Plan du chapitre]
\small
\textbf{I.\quad LA GRAMMAIRE : UN ART OU UNE SCIENCE ?}
\begin{enumerate}[label=\arabic*.]
\item Un art : la grammaire prescriptive
\item Une science : la grammaire descriptive
\item Henri Frei : \textit{La grammaire des fautes}
\end{enumerate}
\textbf{II.\quad GRAMMAIRE ET GRAMMAIRIENS ANTIQUES}
\begin{enumerate}[label=\arabic*.]
\item Les grammairiens indiens : les premiers linguistes ?
\begin{enumerate}[label=\textcircled{\scriptsize\arabic*}]
\item Panini : l’\textit{Ashtadhyayi}
\item Patañjali : le \textit{Mahābhāṣya}
\end{enumerate}
\item Les grammairiens grecs et romains
\begin{enumerate}[label=\textcircled{\scriptsize\arabic*}]
\item Denys le Grammairien
\begin{enumerate}[label=\arabic*.]
\item Élève d’Aristarque et \textit{grammaticus} à Rome
\item Les parties du discours
\end{enumerate}
\item Marcus Varro
\item Aelius Donatus
\begin{enumerate}[label=\arabic*.]
\item Un pionnier de la ponctuation
\item Gutenberg, et l’\textit{Ars Grammatica}
\end{enumerate}
\end{enumerate}
\end{enumerate}
\textbf{III.\quad GRAMMAIRES ET GRAMMAIRIENS DU MOYEN-ÂGE}
\begin{enumerate}[label=\arabic*.]
\item Priscien de Césarée : les \textit{Institutions grammaticales}
\item La grammaire modiste
\begin{enumerate}[label=\textcircled{\scriptsize\arabic*}]
\item Une grammaire universelle
\item Les modes de signification (\textit{modi significandi})
\end{enumerate}
\end{enumerate}
\textbf{IV.\quad GRAMMAIRES ET GRAMMAIRIENS DE LA RENAISSANCE}
\begin{enumerate}[label=\arabic*.]
\item Le français, langue du roi et du droit
\begin{enumerate}[label=\textcircled{\scriptsize\arabic*}]
\item Les premiers édits royaux
\item L’ordonnance de Villers-Cotterêts (1539)
\end{enumerate}
\end{enumerate}
\origpage{59}{045}
\begin{enumerate}[label=\arabic*.,start=2]
\item Grammaires de référence au 16\textsuperscript{e} siècle
\begin{enumerate}[label=\textcircled{\scriptsize\arabic*}]
\item Jacques Dubois : \textit{Une introduction à la langue gauloise}
\item Louis Meigret : \textit{Le tretté de la grammère françoise} (1550)
\begin{enumerate}[label=\arabic*.]
\item Un pionnier de la réforme de l’orthographe
\item Contre l’orthographe latinisante
\end{enumerate}
\item Robert Estienne : \textit{De la précellence du langage françois} (1579)
\item Pierre de La Ramée : \textit{Gramere} (1562)
\end{enumerate}
\item Particularités des grammaires du 16\textsuperscript{e} siècle
\begin{enumerate}[label=\textcircled{\scriptsize\arabic*}]
\item L’article
\item L’adjectif
\item Les pronoms
\item Le participe
\item Prémisses de l’analyse grammaticale
\end{enumerate}
\end{enumerate}
\textbf{V.\quad GRAMMAIRES ET GRAMMAIRIENS AU 17\textsuperscript{e} SIÈCLE}
\begin{enumerate}[label=\arabic*.]
\item Les « censeurs de la langue »
\begin{enumerate}[label=\textcircled{\scriptsize\arabic*}]
\item Claude Favre de Vaugelas : « le greffier de l’usage »
\item Dominique Bouhours : le continuateur de Vaugelas
\end{enumerate}
\item \textit{La grammaire} de Port Royal : le grand tournant
\begin{enumerate}[label=\textcircled{\scriptsize\arabic*}]
\item D’une grammaire « générale et raisonnée » à une grammaire universelle
\item Dualité du signe linguistique
\end{enumerate}
\item L’héritage de Port Royal
\begin{enumerate}[label=\textcircled{\scriptsize\arabic*}]
\item \textit{La linguistique cartésienne} de Noam Chomsky
\item Michel Foucault : Port-Royal et l’« épistémè »
\end{enumerate}
\end{enumerate}
\textbf{Exercices}\par
\textbf{Pour aller plus loin}\par
\textbf{Glossaire}
\end{keybox}
\clearpage
\origpage{60}{046}
\lecturetitle{Introduction}
Le terme « linguistique » n’est apparu qu’au début du 19\textsuperscript{e} siècle. Quant à l’ouvrage considéré comme fondateur de la linguistique moderne, le \textit{Cours de linguistique générale} de Ferdinand de Saussure, il n’a été publié qu’en 1916, il y a à peine plus de 100 ans. Est-ce que cela signifie qu’avant « l’invention » de la linguistique on n’avait jamais fait du langage un objet d’étude ? Bien sûr que non : nous avons vu dans les deux chapitres précédents que depuis toujours, depuis l’« origine », les mythes, les religions, la philosophie, se sont penchés sur le langage. Mais la question du langage n’était pas l’objet spécifique de ces disciplines. Les premiers linguistes, au sens de « spécialistes du langage », ont été les grammairiens antiques de l’Inde, de la Grèce et de Rome. À cette époque, grammaire et linguistique se confondaient, et l’on peut dire que les grammairiens faisaient de la linguistique « sans le savoir », comme Monsieur Jourdain. Au fil du temps, les grammairiens ont pris des chemins différents, et deux conceptions radicalement différentes de la grammaire sont apparues. D’une part, une grammaire normative, faite de règles, qui nous apprend à bien écrire et bien parler une langue particulière. D’autre part, une grammaire descriptive, explicative, dont l’objet est le langage, et qui préfigure notre linguistique moderne.
\section{LA GRAMMAIRE : UN ART OU UNE SCIENCE ?}
Pendant des siècles, un seul terme, celui de « grammaire » a été utilisé pour désigner l’étude de la langue en général et sous tous ses aspects. Aujourd’hui, la grammaire et la linguistique renvoient à deux démarches diamétralement différentes. La première est un art, la seconde est une science.
\subsection{Un art : la grammaire prescriptive}
La définition du terme grammaire a évolué avec le temps. Au 19\textsuperscript{e} siècle, dans \textit{Le Dictionnaire de la langue française} du lexicographe Émile Littré (ainsi appelé aussi le \textit{Littré}), publié entre 1863 et 1872 pour la première édition, on en trouve la définition suivante :
\begin{lecture}{}
L’art d’exprimer ses pensées par la parole ou par l’écriture d’une manière conforme aux règles établies par le bon usage.
\end{lecture}
La grammaire est définie comme un art. Cela implique la possibilité d’un jugement esthétique (beau ou laid) ou éthique (bien ou mal). Par ailleurs, dans cette conception classique, la grammaire est prescriptive, c’est-à-dire qu’elle nous impose des règles à suivre pour « bien » écrire ou « bien » parler sans faire de « fautes ». Autre notion importante dans la définition du \textit{Littré} : le bon usage. Cette grammaire repose sur le parler « correct », le « bon usage » en vigueur pour s’exprimer correctement. Cette grammaire qui a dû vous causer bien des cauchemars lors de votre apprentissage du français n’a \origcont{61}{047}pas de fin scientifique et a pour seul but de dire « comment » on doit s’exprimer. Inutile de s’attarder sur cette grammaire, car c’est celle que nous connaissons et pratiquons tous quotidiennement. C’est elle, par exemple, qui nous apprend le verbe s’accorde avec le sujet,\ednote{原文 « qui nous apprend le verbe s’accorde » 缺连接词，应为 « qui nous apprend que le verbe s’accorde »。} et que « puis-je vous aider » est plus soutenu que « est-ce que je peux vous aider », ou encore que « si l’on » est plus beau que « si on », etc.
\subsection{Une science : la grammaire descriptive}
Faisons à présent un saut dans le temps pour consulter un dictionnaire contemporain, \textit{Le Robert}. Le mot « grammaire » y est défini comme suit :
\begin{lecture}{}
1. Jusqu’au 19\textsuperscript{e} siècle et de nos jours dans le langage courant : ensemble des règles à suivre pour parler et écrire correctement une langue.

2. Étude systématique des éléments constitutifs d’une langue : sons, formes, mots, procédés.
\end{lecture}
On constate qu’à l’époque contemporaine, le mot « grammaire » désigne deux disciplines distinctes. Dans la première partie de la définition, nous retrouvons la grammaire normative (faite de « règles ») et prescriptive (de règles « à suivre ») du Littré. Mais c’est la deuxième définition qui devrait attirer votre attention : la grammaire a une visée non plus prescriptive, mais descriptive. Il s’agit de décrire la langue selon une démarche scientifique (« étude systématique ») et explicative, donc avec rigueur et méthode. Cette grammaire vise l’étude objective de la langue dans ses différentes composantes que sont la phonétique et la phonologie (« les sons »), la morphologie (« les formes »), la sémantique (« les mots »), la syntaxe (« les procédés »), etc.
\subsection{Henri Frei : \textit{La grammaire des fautes}}
Les fautes de grammaire, si elles causent le malheur des étudiants et de leurs professeurs, ont fait l’objet de recherches de la part des linguistes. C’est le cas en particulier d’\textbf{Henri Frei} (1899-1980), linguiste suisse et spécialiste des langues orientales (japonais, hindi, chinois, etc.) qui a publié en 1929 un ouvrage intitulé \textit{La grammaire des fautes}. Dans ce livre, il analyse un corpus important de textes en français populaire constitué de lettres écrites par des soldats de la Première Guerre mondiale. Frei y analyse les fautes non pas du point de vue normatif, comme le ferait un professeur de français corrigeant les copies de ses élèves, mais en tant que linguiste, en adoptant un point de vue descriptif et une méthode scientifique. Dans l’extrait ci-dessous, il développe une approche fonctionnelle de la faute de grammaire :
\begin{lecture}{Lecture : Henri Frei, \textit{La grammaire des fautes} (extrait)}
On ne fait pas des fautes pour le plaisir de faire des fautes. Leur apparition est déterminée, plus ou moins inconsciemment, par les fonctions qu’elles ont à remplir (plus grande expressivité, plus grande clarté, plus grande économie, etc.\ednote{底本在 « remplir » 后打开圆括号，而在 « etc. » 后未闭合；此处应补右括号。引文其余内容保留底本文字。} [...] Dans un grand nombre de cas la faute, qui a passé jusqu’à présent pour un \origcont{62}{048}phénomène quasi-pathologique, sert à prévenir ou à réparer les déficits du langage correct. Partant de faits que le point de vue normatif taxe généralement de fautes, nous chercherons donc, en nous plaçant sur le terrain fonctionnel, à déterminer les fonctions que ces fautes ont à satisfaire. Au contraire, la question de savoir si, et dans quelle mesure, un fait donné est correct ou non, ne nous intéressera que d’une manière secondaire : les ouvrages de purisme abondent, sur ce point.
\end{lecture}
Se plaçant sur le « terrain fonctionnel », il prend ses distances avec « le point de vue normatif » et « les ouvrages de purisme » qui cherchent à savoir « si un fait donné est correct ou non ». Pour Frei, la faute a pour fonction de « prévenir ou à réparer\ednote{本句 « a pour fonction de prévenir ou à réparer » 的并列结构不一致，应作 « a pour fonction de prévenir ou de réparer »。原引文中 « sert à prévenir ou à réparer » 则无此问题。} les déficits du langage correct » : elle contribue donc à faire évoluer la langue. Par exemple, de nos jours, on dit « je me rappelle de cette histoire ». Or, à l’origine, « se rappeler » est un verbe transitif, donc sans préposition : « Je me rappelle cette histoire. » Sous l’influence du verbe « se souvenir », de même sens mais obligatoirement suivi de la préposition « de », l’usage (d’abord fautif) de « se rappeler de », est devenu la nouvelle norme.
\section{GRAMMAIRE ET GRAMMAIRIENS ANTIQUES}
Nous avons défini plus haut le terme de « grammaire ». Dans un instant, son étymologie nous ramènera à la Grèce au temps de Platon. Mais si la Grèce et la Rome antiques ont produit d’illustres grammairiens, le berceau des études sur le langage se trouve en Inde.
\subsection{Les grammairiens indiens : les premiers linguistes ?}
\textbf{Panini} (560-480 avant J.-C.) est sans conteste le grammairien le plus célèbre de l’Inde antique. Il est l’auteur, vers 350 avant J.-C.,\ednote{本段所列生卒年 560–480 avant J.-C. 与“约公元前 350 年”著书不能同时成立：后一年代比所列卒年晚 130 年。此处至少一项有误；具体断代待核。} d’un ouvrage majeur, l’\textit{Ashtadhyayi} (en sanskrit, « composé de huit parties »), composé de 3959 \textit{sutras} (courtes sentences) connues aussi sous le nom de \textit{Paniniya}.
\subsubsection*{1\quad Panini : l’\textit{Ashtadhyayi}}
Dans cet ouvrage, Panini met à jour les règles de morphologie, de syntaxe et de sémantique de la langue de l’Inde antique : le sanskrit. L’\textit{Ashtadhyayi} reste aujourd’hui encore un ouvrage de référence sur la grammaire du sanskrit. Pour beaucoup, Panini serait le premier grammairien ayant jamais existé. Grammaire en huit parties, l’\textit{Ashtadhyayi} est accompagnée de trois autres traités de Panini :
\begin{itemize}
\item Le \textit{Shiva sutra}, qui identifie et liste les phonèmes de la langue sanskrite.
\item Le \textit{Dhatu paṭha}, qui contient les racines verbales du sanskrit.
\item Le \textit{Gaṇapaṭha}, qui s’intéresse à la dérivation (formation de nouveaux mots) en sanskrit.
\end{itemize}
L’\textit{Ashtadhyayi} est concis mais difficile à aborder. Il a été rendu compréhensible grâce à deux commentaires :
\begin{itemize}
\item Le \textit{Vārttika}, écrit par le grammairien \textbf{Kātyāyana}.
\item Le \textit{Mahābhāṣya} (\textit{le Grand Commentaire}) écrit par le grand savant indien \textbf{Patañjali}.
\end{itemize}
\origpage{63}{049}
La grammaire de Panini est fortement systématisée et technique. Les concepts de phonème, de morphème, de racine et de flexion, ne seront pris en considération par les linguistes occidentaux que deux millénaires après lui. Les règles définies par Panini sont si claires que les informaticiens les utilisent de nos jours en intelligence artificielle pour enseigner la compréhension du sanskrit aux ordinateurs. Panini emploie des métarègles, des transformations et la récursivité ainsi qu’un système élaboré d’abréviations. La forme de Backus-Naur, système de notation développé par \textbf{John Backus} et \textbf{Peter Naur} dans les années 50 pour décrire les règles syntaxiques des langages de programmation et qui est utilisé dans de nombreux logiciels d’analyse syntaxique, possède des similitudes importantes avec les règles de la grammaire de Panini. Pour cette raison, Panini est considéré comme un des pères de l’informatique.
\subsubsection*{2\quad Patañjali : le \textit{Mahābhāṣya}}
Le \textit{Mahābhāṣya} de \textbf{Patañjali} postule que les mots sont indépendants de la phonation, c’est-à-dire des émissions sonores qui les expriment. Cette distinction n’est pas sans rappeler la distinction entre « signifiant » et « signifié » chez Ferdinand de Saussure, qui était justement professeur de sanskrit. Les théories de Patañjali sur le langage et la parole offrent de multiples résonances avec les théories linguistiques modernes, en particulier la théorie des universaux du langage.
\subsection{Les grammairiens grecs et romains}
Si la grammaire est née en Inde, c’est de Grèce (et de Platon) que nous vient le terme de « grammaire ». La \textit{grammatikè technè} en Grec, c’est littéralement « l’art (\textit{technè}) des lettres (\textit{grammata}) ». La grammaire était à l’origine un art de l’écriture et de la lecture faisant l’objet d’un enseignement de professeurs appelés \textit{grammatistès}.

L’étymologie du mot grammaire met en évidence une caractéristique importante de cet art, et cela dès sa naissance : son objet par excellence est l’écrit. Par conséquent, la naissance de la grammaire est contemporaine de l’introduction de l’écriture chez les Grecs. Les Grecs ont emprunté leur écriture aux Phéniciens dans les premiers siècles du Ier millénaire avant notre ère, car ils avaient perdu leur propre écriture, le mycénien. Perdu pendant de très nombreux siècles, le mycénien a été redécouvert en 1900 par \textbf{Sir Arthur Evans} sur des tablettes d’argile dans la ville de Cnossos, lieu de naissance du mycénien.
\subsubsection*{1\quad Denys le Grammairien}
\textbf{Denys le Grammairien} (170-90 avant J.-C.) est un illustre grammairien grec originaire de Thrace. Il est né à Alexandrie, en Égypte, et cela n’est pas sans conséquence.
\minititle{1. Élève d’Aristarque et \textit{grammaticus} à Rome}
Denys a été le disciple d’\textbf{Aristarque de Samothrace}, (217-144 avant J.-C.), grammairien, philologue et pionnier de l’étude des textes d’Homère. Bénéficiant de l’enseignement d’un tel grand maître, Denys a exercé le métier de \textit{grammaticus} à Rome.
En tant que \textit{grammaticus}, il enseignait aux adolescents la langue écrite, la prononciation et la diction du grec et du latin et \origcont{64}{050}leur faisait connaître les grands poètes grecs (Homère) et romains (Virgile), à travers l’analyse et l’explication de textes.\ednote{此处将不同年代的教学情形混在一起。原文所列 Denys 卒年为公元前 90 年，而 Virgile 生于公元前 70 年（Biblissima 的作者规范记录，见\ecite{14}），晚于前述卒年；因此不能把讲授 Virgile 作品写作 Denys 本人的经历。应将一般的罗马语法教学概述与 Denys 的传记分开。} À l’époque de la république romaine, un seul enseignant enseignait deux langues mais sous l’Empire, ils se sont spécialisés : on trouvait des \textit{grammaticus Graecus} (professeurs de grec) et des \textit{grammaticus Latinus} (professeurs de grec).\ednote{第二个 « professeurs de grec » 显为重复误写；与 « Latinus » 对应，应作 « professeurs de latin »。} À cette époque, le grec était la langue des élites romaines. C’était également le cas au 17\textsuperscript{e} siècle en France.
\minititle{2. Les parties du discours}
On doit à Denys une grammaire de la langue grecque longtemps classique et qui a servi de manuel de référence dans l’étude du grec. Mais au-delà de son apport à la grammaire normative, un des apports directs de Denys à la linguistique est l’identification des parties du discours. Selon lui, tout mot relève de l’une de ces huit catégories : nom, verbe, participe, article, pronom, préposition, adverbe, conjonction. Sa liste des parties du discours a été utilisée jusqu’en 1780. Cette année-là, dans ses \textit{Éléments de la grammaire françoise}, le grammairien \textbf{Charles François Lhomond} (1727-1794) a adapté les catégories de Denys à la langue française, et a proposé neuf classes, ajoutant l’interjection. La liste des parties du discours en français est ensuite devenue canonique en 1910, avec le retrait du participe et l’ajout de l’article.\ednote{此处叙述内部不一致：本段最初八类已经包含 « article »，后面却再把 « l’ajout de l’article » 作为新变化。不能据此直接重建各时期完整词类表；有关 Lhomond 及 1910 年分类的具体对应仍待版本核验。}
\subsubsection*{2\quad Marcus Varro}
Savant, homme politique et écrivain prolifique (il a écrit plus de 600 ouvrages), \textbf{Varro} est l’auteur d’une grammaire du latin, \textit{De lingua latina} (« De la langue latine »), qui a longtemps été une référence incontournable pour les grammairiens latins, même si seulement le quart est parvenu jusqu’à nous. Varro y distingue pour la première fois les mots variables et invariables :
\begin{lecture}{Lecture : \textit{De lingua latina} (extrait)}
On conçoit, en effet, que les mots servant à désigner des idées invariables devaient être également invariables, de même que, dans une maison où il n’y a qu’un seul esclave, cet esclave n’a besoin que d’un nom ; tandis que, dans une maison où il y en a plusieurs, chaque esclave a besoin de plusieurs noms, pour qu’on puisse le distinguer de ses compagnons. Ainsi les mots et les noms qui expriment des idées variables doivent nécessairement subir des modifications correspondantes à ces idées ; tandis que les mots qui ne servent qu’à unir les mots entre eux, sont ordinairement invariables et ressemblent à une courroie, qui peut également servir à attacher un homme, un cheval, etc. Quand nous disons, par exemple : sous le consulat de Tullius et d’Antonius, nous sentons que la conjonction « et » peut unir non seulement les noms de deux consuls quelconques, mais encore tous les noms et tous les mots sans exception.
\end{lecture}
\subsubsection*{3\quad Aelius Donatus}
\textit{Grammaticus} romain, Donat (320-380) a été le professeur de \textbf{Jérôme de Stridon} (347-420), le traducteur de la \textit{Bible} en latin (\textit{la Vulgate}), plus connu sous le nom de Saint Jérôme. Mais il est surtout l’auteur de l’\textit{Ars grammatica}, une grammaire latine tellement célèbre que le nom propre \origcont{65}{051}« Donat » est devenu un nom commun, « un donat », synonyme de livre de grammaire. C’est pour cela que la plus ancienne grammaire du français est connue sous le nom de « Donait françois » (1409), c’est-à-dire le « donat français ».
\minititle{1. Un pionnier de la ponctuation}
Donat (ou Donatus) est l’un des premiers à avoir utilisé un système archaïque de ponctuation. À partir de la fin du 1er siècle de notre ère, l’écriture était sans blancs ni espaces : c’est la \textit{scriptio continua} (« effet continue »).\ednote{原文 « effet continue » 既未正确释出 scriptio，形容词亦不与阳性名词 effet 配合。此拉丁术语在这里应释为 « écriture continue »，即不以空格分词的连续书写。} Aux 3\textsuperscript{e} et 2\textsuperscript{e} siècles avant J.-C., \textbf{Aristophane de Byzance} (257-180 avant J.-C.) et Aristarque de Samothrace, successivement directeurs de la Bibliothèque d’Alexandrie, avaient déjà défini un système comportant trois types de points pour marquer la ponctuation : le point d’en haut pour la fin d’une phrase, le point médian marquant une pause moyenne et le point d’en bas pour marquer une courte pause. Donat a été l’un des premiers utilisateurs de ce système de ponctuation dans l’alphabet romain : un point en haut pour une pause forte, un point au milieu de la ligne pour une pause médiane et un point sur la ligne pour une pause faible : cela correspond aujourd’hui à la virgule, au point-virgule et au point.\ednote{若按本句先列的“强—中—弱停顿”顺序类比现代标点，末尾应按 « au point, au point-virgule et à la virgule » 排列，而非原文的相反顺序；古今标点功能也不宜理解为完全一一对应。}

Ce système a été complété par des lettres d’annotation : trois lettres en particulier, K, C et $\Gamma$ (\textit{gamma}) signalaient le début d’un paragraphe, c’est-à-dire une unité textuelle de plus grande envergure que la phrase. La lettre elle-même était placée dans le texte ou en marge de la ligne où commençait le paragraphe, mais sans retour à la ligne, comme c’est le cas de nos jours. L’espace blanc et le retour à la ligne que nous connaissons aujourd’hui n’interviendront que plus tard comme moyens de ponctuation. Le système de Donat a été utilisé jusqu’au 7\textsuperscript{e} siècle, où il a été remplacé par un système plus performant, celui d’\textbf{Isidore de Séville} (560-636) dans le haut-moyen âge.
\minititle{2. Gutenberg, et l’\textit{Ars Grammatica}}
Tout le monde sait que le premier livre imprimé entre 1452 et 1455 par \textbf{Johannes Gutenberg} (1400-1468) était une \textit{Bible}. La raison était matérielle : Gutenberg devait imprimer un livre dont le tirage (et donc les ventes) lui permettrait de couvrir les sommes engagées. À l’époque, le seul livre capable d’un succès immédiat était la \textit{Bible} de Saint Jérôme. Il faut savoir qu’avant l’invention par Gutenberg de l’impression, un moine copiste avait besoin d’environ trois ans pour recopier à la main une \textit{Bible}. Mais en réalité, c’est un donat, c’est-à-dire la grammaire latine de Donat, qui a été le tout premier livre imprimé par Johannes Gutenberg en 1451 pour tester son invention avant de se lancer dans l’impression de la \textit{Bible}.
Cela devrait vous montrer l’importance de ce manuel dans l’Europe du Moyen Âge et de la Renaissance.
\section{GRAMMAIRES ET GRAMMAIRIENS DU MOYEN ÂGE}
Toutes les grammaires imprimées au Moyen Âge se sont inspirées de l’\textit{Ars grammatica} de Donat, et c’est le cas des \textit{Institutiones grammaticae} de Priscien de Césarée, grammaire de référence dans toute l’Europe médiévale.
\origpage{66}{052}
\subsection{Priscien de Césarée : les \textit{Institutions grammaticales}}
\textbf{Priscien de Césarée} est un grammairien latin du 6\textsuperscript{e} siècle. Suite au sac de Rome par les Vandales en 455, il a été exilé et s’est installé à Constantinople. Il y a ouvert en 525 une école latine publique renommée. Son principal ouvrage est sa grammaire, les \textit{Institutiones grammaticae} (« institutions grammaticales »), ouvrage très célèbre dont les copies manuscrites se sont diffusées dans toute l’Europe. On dispose aujourd’hui de 61 manuscrits complets ou fragmentaires écrits entre 790 et 900. Quant aux plus récents, ils se comptent par centaines, ce qui atteste de la très grande popularité de l’ouvrage. Les seize premiers chapitres des \textit{Institutiones grammaticae} traitent des lettres, des syllabes, des parties du discours. Les deux derniers, intitulés « De Constructione », traitent de la syntaxe.
\subsection{La grammaire modiste}
Au Moyen Âge, avec la redécouverte des œuvres Aristote\ednote{原文缺少表示作者关系的介词，应为 « des œuvres d’Aristote »。} et les progrès subséquents dans le domaine de la logique, une forme nouvelle de grammaire est apparue : une grammaire dont l’objet n’est pas une langue particulière, mais le langage en général.
\subsubsection*{1\quad Une grammaire universelle}
Les grammairiens vont commencer à distinguer le langage, qui est un phénomène « universel », et la langue, variation « accidentelle », pour reprendre la terminologie d’Aristote. Enseignée durant la seconde moitié du 13\textsuperscript{e} siècle et la première moitié du 14\textsuperscript{e} siècle à Paris, la grammaire modiste (appelée aussi « spéculative ») se proposait de fonctionner comme un miroir donnant l’image vraie du langage au-delà de la simple description des faits langagiers particuliers directement perceptibles. Pour \textbf{Robert Kilwardby} (1215-1279) archevêque de Cantorbéry, cardinal, professeur de grammaire et de logique, la grammaire modiste est une grammaire s’abstrayant de toute langue particulière. De même, le grand savant anglais \textbf{Roger Bacon} (1214-1294), la décrit comme « une grammaire unique et la même chez tous », et pose de façon implicite son universalité :
\begin{lecture}{}
La grammaire est substantiellement la même dans toutes les langues, bien qu’elle y subisse des variations accidentelles.
\end{lecture}
Cette intuition très précoce aura, des siècles plus tard, et de nos jours encore, une influence considérable en linguistique. Les grammairiens modistes étudiaient donc les traits communs à toutes les langues, laissant à la grammaire normative les règles particulières à chaque langue. En essayant d’expliquer les fondements de la grammaire et en énonçant des règles de portée générale, les modistes ont donc contribué à transformer un art en science. En tant que science, il leur fallait en effet des principes, un objet, des méthodes propres. Le fonctionnement de la grammaire modiste reposait en particulier sur l’articulation sur deux niveaux : dans le premier, il définissait ses unités minimales, c’est-à-dire les parties du discours ; dans le second, il établissait les règles de construction des énoncés à partir des unités minimales. C’est ce qui correspond de nos jours à la syntaxe, que l’on peut définir comme l’ensemble des règles régissant la combinaison des unités linguistiques pour former des phrases. Les grammaires génératives contemporaines, en particulier \origcont{67}{053}la \textit{generative grammar} de \textbf{Noam Chomsky}, ont été fortement influencées par la grammaire des modistes.
\subsubsection*{3\quad Les modes de signification (\textit{modi significandi})}
\ednote{正文此处的带圈序号在底本作 3，而本节此前只有第 1 小项；章首提纲相应条目作 2。序号应为 2，此处保留底本正文的 3 并加注说明。}
Pour les modistes, les principes de la grammaire sont appelés \textit{modi significandi}, c’est-à-dire \textit{modes de signification}. C’est de ces « modes » que nous vient le nom de « grammaire modiste ». Selon les modistes, les objets du monde ont des propriétés qui peuvent être envisagées sous trois modes :
\begin{itemize}
\item Les propriétés essentielles, qui relèvent du mode d’être (\textit{modus essendi}). Il s’agit de leur essence.
\item Les propriétés appréhendées par notre esprit. Elles relèvent du mode de compréhension (\textit{modus intelligendi}).
\item Les propriétés matérialisées par les mots en parties du discours et catégories grammaticales. Elles relèvent du mode de signification (\textit{modus significandi}).
\end{itemize}
Par exemple, « souffrir », « souffrant », « souffrance » renvoient à la même notion mais se manifestent par différentes parties du discours (respectivement verbe, adjectif, et nom). Ils ont donc des modes de signification distincts. Ces propriétés ne sont pas exclusivement rattachées à un objet extérieur : elles peuvent être étudiées pour elles-mêmes. En d’autres termes, la grammaire modiste remplace l’étude de la signification, qui s’appuie sur une réalité extérieure, par l’étude des modes de signification qui, eux, relèvent du fonctionnement interne de la langue. Parmi les grammairiens modistes les plus illustres, \textbf{Martin de Dacie}, auteur de \textit{De modis significandi} (vers 1255) est souvent considéré comme le premier des modistes. Parmi les modistes français, citons deux professeurs à l’Université de Paris : \textbf{Raoul le Breton} (1230-1320), et \textbf{Siger de Courtrai} (1283-1341), auteur d’une \textit{Summa modorum significandi} (« Somme sur les modes de signification »).

Si le Moyen Âge du 13\textsuperscript{e} et 14\textsuperscript{e} siècle a vu l’apparition de cette grammaire d’un nouveau type, le 15\textsuperscript{e} siècle a vu la publication de la toute première grammaire du français.
\begin{bbcomplement}{}
Le \textit{Donait françois} (1409) est unanimement reconnu comme la plus ancienne grammaire du français, écrite par un anglais (\textbf{John Barton}) pour des Anglais (« les bones gens du roiaume d’Engleterre »). En voici la préface :
\minititle{John Barton, \textit{Donat français} (extrait)}
« Pour ceo que les bones gens du roiaume d’Engleterre sont embrasez a sçavoir lire et escrire, entendre et parler droit françois, afin qu’ils puissent entrecomuner bonement ové lour voisins, c’est a dire les bones gens du roiaume de France, et ainsi pour ce que les leys d’Engleterre pour le graigneur partie et aussi beaucoup de bones choses sont misiez en françois, et aussi bien pres touz les seigneurs toutes les dames \origcont{68}{054}en mesme roiaume d’Engliterre volentiers s’entrescrivent en romance, tres necessaire je cuide estre aus Engleis de sçavoir la droite nature de françois. Au honneur de Dieu et de sa tres doulce miere et toutz les saintez de paradis, je, Johan Barton, escolier de Paris, nee et nourie toutez voiez d’Engleterre en le conté de Cestre, j’ey baillé aus avant diz Anglois un Donait françois pur les briefment entreduyr en la droit language du Paris et de païs la d’entour, la quelle language en Engliterre on appelle doulce France. »

Cette préface est écrite en anglo-normand, ancienne langue d’oïl (c’est-à-dire du Nord de la France) parlée au Moyen Âge en Angleterre à la Cour des rois et dans l’aristocratie anglo-normande. La conquête du royaume d’Angleterre en 1066 par le duc de Normandie, Guillaume le Conquérant, avait eu pour conséquence l’introduction en Angleterre du langage normand. C’est ce « normand insulaire » qu’on appelle anglo-normand.
\end{bbcomplement}
\section{GRAMMAIRES ET GRAMMAIRIENS DE LA RENAISSANCE}
\markright{Grammairiens de la Renaissance}
Au tout début de la renaissance, le \textit{Donat français} de Barton a fait des émules, et la réflexion grammaticale a été favorisée par plusieurs éléments. D’une part, les efforts d’unification linguistique menés par François Ier (1494-1547) pour que le français devienne la langue administrative et la langue du pouvoir centralisé. D’autre part, l’affirmation d’une littérature française, avec le manifeste de \textbf{Du Bellay} (1522-1560), \textit{Défense et illustration de la langue française} (1549).
\subsection{Le français, langue du roi et du droit}
A la Renaissance, le pouvoir royal s’intéresse au français. Comme l’Église avec le latin universel, la monarchie française a elle aussi choisi sa langue véhiculaire pour uniformiser juridiquement et administrativement le pays. Le français est ainsi devenu la langue du droit et du roi.
\subsubsection*{1\quad Les premiers édits royaux}
En décembre 1490, l’ordonnance de Moulins du roi Charles VIII (1470-1498) obligeait l’emploi du « langage françois ou maternel » et non le latin dans les interrogatoires et procès-verbaux dans la région du Languedoc. En 1510, le roi Louis XII (1462-1515) avait lui aussi exigé l’emploi du « vulgaire et langage du pays » :
\begin{lecture}{}
Ordonnons que doresnavant tous les procez criminels et lesdites enquestes, en quelque maniere que ce soit, seront faites en vulgaire et langage du pais... autrement ne seront d’aucun effet ni valeur.
\end{lecture}
En 1535, l’ordonnance d’Is-sur-Tille de François I\textsuperscript{er} (1494-1547) prescrivait que les actes judiciaires et administratifs soient rédigés « en françoys ou à tout le moins en vulgaire dudict pays ». Cette mesure royale faisait en sorte que les procédures judiciaires et les décisions de justice soient accessibles à la population. Pour cela, il fallait utiliser le « langage maternel francoys » au lieu du latin. Il faut savoir qu’à l’époque, le français était aussi étrange que le latin pour l’immense
\origcont{69}{055}majorité de la population, et la plupart des ordonnances royales précédentes (entre 1490 et 1535) autorisaient le choix entre deux usages linguistiques :
\begin{itemize}
\item « en langage François ou maternel » (ordonnance de 1490) ;
\item « en vulgaire ou langage du païs » (ordonnance de 1510) ;
\item « en langue vulgaire des contractants » (ordonnance de 1531) ;
\item « en francoys ou a tout le moins en vulgaire dudict pays » (ordonnance de 1535).
\end{itemize}
Notons que l’ajout de la mention « à tout le moins » dans l’ordonnance d’Is-sur-Tille marque une hiérarchie dans l’esprit du roi entre le français et les parlers régionaux.
\subsubsection*{2\quad L’ordonnance de Villers-Cotterêts (1539)}
On ignore encore aujourd’hui si le « langage maternel francoys » désignait la langue maternelle du roi, celle de la population de l’Île-de-France, ou celle de tous les Français. À l’époque les patois (parlers locaux) étaient omniprésents.

C’est l’ordonnance royale de Villers-Cotterêts (1539) qui a le plus marqué le statut du français, et c’est dans son château de Villers-Cotterêts que François I\textsuperscript{er} (qui parlait le françoys, le latin, l’italien et l’espagnol) l’a signée, imposant ainsi « le françoys » comme langue administrative au lieu du latin. L’ordonnance royale de Villers-Cotterêts traite certes de la langue, mais de façon très marginale : seuls 2 articles sur les 192 qu’elle contient. Le titre de l’ordonnance mentionne clairement qu’elle porte en fait sur la justice : \textit{Ordonnance du Roy sur le faict de justice}.
\begin{lecture}{Lecture : \textit{Ordonnance de Villers-Cotterêts} (1539), articles 110 et 111}
\textnormal{\textbf{Version originale}}\par
110 Que les arretz soient clers et entendibles et afin qu’il n’y ayt cause de doubter sur l’intelligence desdictz Arretz, nous voullons et ordonnons qu’ilz soient faictz et escriptz si clerement qu’il n’y ayt ne puisse avoir aulcune ambiguite ou incertitude, ne lieu a en demander interpretacion.

111 Nous voulons que doresenavant tous arretz, ensemble toutes aultres procedures, soient de noz courtz souveraines ou aultres subalternes et inferieures, soient de registres, enquestes, contractz, commissions, sentences, testamens et aultres quelzconques actes et exploictz de justice ou qui en deppendent, soient prononcez, enregistrez et delivrez aux parties en langaige maternel francoys et non aultrement.

\textnormal{\textbf{Translittération}}\par
110 Afin qu’il n’y ait cause de douter sur l’intelligence des arrêts de nos cours souveraines, nous voulons et ordonnons qu’ils soient faits et écrits si clairement, qu’il n’y ait ni puisse avoir ambiguïté ou incertitude, ni lieu à demander interprétation.

111 Nous voulons donc que dorénavant tous arrêts, et ensemble toutes autres procédures, soient de nos cours \origcont{70}{056}souveraines ou autres subalternes et inférieures, soient des registres, enquêtes, contrats, testaments et autres quelconques actes et exploits de justice ou qui en dépendent, soient prononcés, enregistrés et délivrés aux parties en langage maternel françoys et non autrement.\ednote{底本所标 « Translittération » 实为现代化转写，且第 111 条漏掉前列原文中的 « commissions, sentences »。与本页原文对齐时，« contrats » 后应补 « commissions, sentences »；这不是据外部版本改动旧法文。}
\end{lecture}
\subsection{Grammaires de référence au 16\textsuperscript{e} siècle}
Parmi le très grand nombre de grammaires du français publiées au 16\textsuperscript{e} siècle, on en présentera quatre qui ont fait date.
\subsubsection*{1\quad Jacques Dubois : \textit{Une introduction à la langue gauloise}}
Si les grammaires de John Barton (auteur de \textit{Donat français}) et John Palsgrave (auteur de \textit{Lesclarcissement de la langue francoyse}) ont été écrites en anglais et pour des Anglais,\ednote{将两书一并说成 « écrites en anglais » 与底本前两页展示的 Barton 前言冲突：该段文字及作者自己的说明均将其认作 anglo-normand，而非英语。应区分“面向英国读者”与“用英语写成”。} de nombreux grammairiens français continuaient quant à eux à écrire leurs grammaires du « françois » en latin : c’est le cas, par exemple, de Jacques Dubois.

Médecin, anatomiste français enseignant à la Faculté de médecine de Paris, \textbf{Jacques Dubois} (1478-1555), a étudié les arts libéraux au Collège de Tournai. En plus du latin qu’il maîtrisait parfaitement, il a appris le grec et l’hébreu. Il est aussi l’auteur de la première grammaire du français écrite en France par un Français et publiée en 1531 par le célèbre imprimeur \textbf{Robert Estienne} (1503-1559) : \textit{Une introduction à la langue gauloise ainsi que sa grammaire à partir d’auteurs hébreux, grecs et latins}. Comme son titre l’indique, il s’agit d’une d’étude\ednote{原文 « d’une d’étude » 多出第二个 d’，应为 « d’une étude »。} grammaticale de la langue française dans laquelle il s’appuie principalement sur la comparaison avec les langues antiques pour expliquer les spécificités du français. À cette époque, on ne concevait pas encore de décrire une langue moderne autrement que par rapport aux modèles grec et latin. Et c’est d’ailleurs en latin qu’il a écrit sa grammaire du français, « pour que ces principes de notre langue puissent servir à la fois aux Anglais, aux Allemands, aux Italiens, aux Espagnols, à tous les étrangers enfin ». Dans sa grammaire, Dubois est le premier à avoir utilisé l’apostrophe, le tréma et l’accent circonflexe.
\subsubsection*{2\quad Louis Meigret : \textit{Le tretté de la grammère françoise} (1550)}
\textbf{Louis Meigret} (1510-1558) a écrit en 1550 la première grande grammaire du français en français sous le titre \textit{Le tretté de la grammère françoise}.
\minititle{1. Un pionnier de la réforme de l’orthographe}
Défenseur de la réforme de l’orthographe française, Meigret a cherché à la simplifier en proposant une orthographe phonétique. Rencontrant l’opposition farouche des poètes de la Pléiade \textbf{Jacques Peletier du Mans} (1517-1582) ainsi que celle de \textbf{Guillaume Des Autels} (1529-1580), la réforme de l’orthographe de Meigret n’a pu s’imposer, faute d’autorité linguistique au 16\textsuperscript{e} siècle. L’orthographe est donc restée étymologique, et certains imprimeurs de l’époque l’ont compliquée un peu en utilisant non seulement du latin, mais également du grec.
\origpage{71}{057}
\minititle{2. Contre l’orthographe latinisante}
Si Meigret a voulu simplifier l’orthographe, c’est parce que jusqu’alors, tous les clercs, érudits, juristes, copistes, poètes et autres écrivaient les mots comme bon leur semblait, sans aucune contrainte. Par ailleurs, les imprimeurs, on vient de le voir, introduisaient des « consonnes étymologiques » dans l’orthographe des mots français, alors qu’elles n’étaient pas prononcées. Par exemple, ils ajoutaient un « g » et un « t » dans le mot « doi » (doigt) pour rappeler son origine latine (\textit{digitum}). De même, pour :
\begin{itemize}
\item le « p » de compter (\textit{computare}) ;
\item le « g » de « cognoistre » (\textit{cognoscere}) ;
\item le « p » de « corps » (\textit{corpus}) ou de « temps » (\textit{tempus}) ;
\item le « h » de « homme » (\textit{homo}), le « b » de « soubdain » (\textit{subitaneus}) ;
\item le « p » de « sept » (\textit{septem}), le « g » de « vingt » (\textit{viginti}) ;
\item le « x » de « paix » (\textit{pax}), etc.
\end{itemize}
En ancien français, ces mots s’écrivaient respectivement : « doi », « conter », « conoitre », « cors », « tems », « om », « sudein », « set », « vint » et « pais ». En ajoutant des consonnes étymologiques, les latiniseurs multipliaient donc dans l’orthographe les lettres superflues rappelant l’étymologie latine et grecque. Ces complications orthographiques étaient une bonne affaire pour les imprimeurs et les typographes qui pouvaient afficher leurs connaissances et surtout gagner plus d’argent, puisque les mots d’imprimerie étaient plus longs à composer. À cette époque, les écrivains, les lettrés et les gens du monde ont cédé en faveur des typographes et leur ont laissé le soin d’écrire le français comme ils le jugeaient, c’est-à-dire de façon plus savante, plus latinisante, plus noble. L’\textit{Apologie de Raymond Sebond} de Montaigne nous offre de très nombreux exemples de cette écriture : « nostre maladie » (\textit{noster}) ; « traictent » (\textit{tractatus}); « cognoist » (\textit{cognoscere}) ; « void » (\textit{video}), etc.
\subsubsection*{3\quad Robert Estienne : \textit{De la précellence du langage françois} (1579)}
\label{ch03-estienne}
\textbf{Robert Estienne} (1503-1559) était le fils de l’imprimeur et helléniste Henri Estienne. Imprimeur royal de François Ier, il était spécialisé dans les ouvrages en hébreu, latin, et grec. Imprimeur, on l’a vu, de la grammaire de Jacques Dubois, il ne s’est pas contenté d’éditer les ouvrages des autres et a apporté son propre \textit{Traicté de la grammaire Françoise}, publié en 1569. Dans \textit{De la précellence du langage françois} (1579), Robert Estienne\ednote{本小节混合了不同作者。1579 年《De la précellence du langage françois》的作者为 Henri Estienne，而非此处所述 1503–1559 年的 Robert Estienne；所据 1579 年版本书名及书目记录明确署名 Henri Estienne，见\ecite{13}。1569 年语法书的版本归属另需分别核对，不能把本段所有 « Robert » 一律改成 « Henri »。} estime que les patois constituent une richesse mais que le « françois » doit demeurer la langue principale du royaume :
\begin{lecture}{}
J’estime qu’en cas de langage je ne puis appeler le cœur de la France les lieux où sa nayveté et pureté est le mieux conservee : de sorte que tous y sont d’accord que ces voscables estrangers nous doivent servir de passetemps plustost que d’ornement ou enrichissement, et que le langage de ceux qui en usent autrement, doit estre déclaré non pas françois mais gastefrançois.
\end{lecture}
\origpage{72}{058}
En 1530, pour le malheur des étudiants en français, c’est Robert Estienne qui a introduit en typographie l’accent aigu, l’accent grave et a perfectionné l’accent circonflexe.
\subsubsection*{4\quad Pierre de La Ramée : \textit{Gramere} (1562)}
Logicien et philosophe français, \textbf{Pierre de La Ramée} (1515-1572) est l’un des plus grands savants humanistes du 16\textsuperscript{e} siècle. Il est aussi le premier à avoir proposé de distinguer le « u » du « v » et le « i » du « j », proposition qui figure dans sa \textit{Gramere} parue en 1562. Il a cependant fallu attendre 1718 pour que sa réforme soit adoptée. Les lettres « j » et « v », qui n’existaient pas dans l’alphabet latin, sont depuis appelées consonnes « ramistes ».

Une partie importante de l’œuvre de La Ramée, ce sont ses grammaires : une grammaire latine publiée en 1559, une grammaire grecque publiée en 1560, et une grammaire française, publiée en 1562 sous le titre de Gramere. Dix ans plus tard, elle a été rééditée avec une autre orthographe : \textit{Grammaire de P. de La Ramée, lecteur du Roy en l’Université de Paris}. Dans cet ouvrage, La Ramée poursuit les tentatives de ses prédécesseurs (Dubois, Meigret, Pelletier) de régenter la langue et l’orthographe françaises. L’édition de 1572 présente une impression sur deux colonnes juxtaposant l’orthographe ramiste et l’orthographe traditionnelle. Sous la forme d’un dialogue entre un précepteur et un disciple, ce traité comporte deux livres (le premier consacré à la morphologie, le second à la syntaxe) et expose de façon méthodique les lettres, les syllabes, les conjugaisons, les articles, les accords, le genre et le nombre des noms, les pronoms, et les prépositions.
\subsection{Particularités des grammaires du 16\textsuperscript{e} siècle}
De manière générale, le 16\textsuperscript{e} siècle adopte le modèle des grammaires latines inspirées de \textit{Donat}, tout en l’adaptant. À cette époque, le statut de l’article est encore instable, l’adjectif est absent et le participe est considéré comme une partie du discours à part entière. Rappelons que la liste des parties du discours en français n’est devenue canonique qu’en 1910, avec le retrait du participe et l’ajout de l’article.
\subsubsection*{1\quad L’article}
Un article est une partie du discours qui précède le substantif (ou l’adjectif précédant le substantif) et en précise le genre (masculin, féminin), le nombre (singulier, pluriel), et la notoriété (article défini, article indéfini).

La liste des parties du discours dans les grammaires latines ne contenait pas l’article, puisqu’il n’y en a pas en latin. Cela avait constitué un problème pour les grammairiens du Moyen Âge, et cela était encore le cas au 16\textsuperscript{e} siècle. Parmi les auteurs que nous avons présentés plus haut, seul Estienne reconnaît à l’article le statut de partie du discours. Les autres grammairiens le traitent comme une particule accompagnant le nom et l’intègrent dans cette catégorie. Mais lorsque l’on se penche sur ce que les grammairiens de l’époque désignent par « article », on ne peut qu’être très surpris. En effet, pour certains auteurs, comme Sylvius et Estienne, l’article prend des formes variées : « le » maistre (maître), « de » maistre, « du » maistre, « au » maistre, « a » maistre, « o » maistre, etc. Pour Le Ramée,\ednote{此处人名 « Le Ramée » 与本节及章首的 « La Ramée » 不一致，应为 « La Ramée »。} ces déterminants contractés (« du », « des » et « aux ») sont des « prépositions » s’utilisant directement devant un nom, qui est alors dépourvu d’article.
\origpage{73}{059}
\subsubsection*{2\quad L’adjectif}
Un adjectif est une partie du discours variable en genre et en nombre, l’adjectif se rapporte au substantif (nom) et en désigne la qualité.

Si l’on cherche dans la liste des parties du discours que proposent les grammaires du 16\textsuperscript{e} siècle, on n’y trouvera pas l’adjectif. En fait, il est inclus dans la section consacrée aux noms. Cette classe, très vaste à l’époque, englobait les « noms substantifs » et les « noms adjectifs » :
\begin{lecture}{Les adjectifs – Estienne, \textit{Traicté De la grammaire Françoise} (1569)}
Il y a deux sortes de noms : les uns sont appelés substantifs, desquels la signification est entendue sans qu’autres mots leur soient adjoints, comme pain, terre, et font un sens parfait avec l’adjectif, comme pain blanc, terre noire. Les adjectifs sont les mots qui se mettent avec les substantifs pour déclarer leur qualité ou quantité, et ne se mettent point proprement sans substantif ou autrement on ne saurait à quoi servirait ledit adjectif, comme en disant blanc, tu ne peux rien entendre si tu n’adjoins quelque substantif, comme disant pain blanc, terre noire, terre grasse, bon homme, homme juste, mauvaise personne, grand personnage, grand larron, vin excellent, homme prudent, riche, pauvre, et ainsi des autres.
\end{lecture}
\subsubsection*{3\quad Les pronoms}
Un pronom est une partie du discours ayant la propriété de remplir les mêmes fonctions que le nom ou le syntagme nominal qu’il remplace dans une phrase. Il peut soit désigner directement quelqu’un ou quelque chose (pronom déictique), soit représenter un segment du discours (pronom anaphorique).

Pour un lecteur moderne, l’analyse des pronoms faite par les grammairiens du 16e siècle est elle aussi surprenante. La logique de classement repose essentiellement sur le principe selon lequel, selon Estienne, les pronoms « dénotent toujours quelque certaine personne ». La liste contient ainsi les pronoms de base (« primitifs »), qui représentent la première, la deuxième et la troisième personne de la conjugaison (je, tu, il, elle, etc.), mais aussi, et c’est plus surprenant, tout mot formé sur ces personnes. Nos adjectifs possessifs modernes (mon, ton, son, sien, etc.) faisaient donc partie à l’époque des pronoms, puisque leur formation reposait entre autres sur la personne à laquelle ils faisaient référence.
\begin{lecture}{Les pronoms – Louis Meigret, \textit{Le tretté de grammere françoèze} (1550)}
Sont dérivés de la première personne, mon, ma, de moi, ou me, et mien, mienne ; et de nous, nos, notre. Et de la seconde personne tu, ou toi, ou te : ton, ta, tien ; et de vous, son pluriel : vos, votre. Reste le réciproque soi, qui fait son, sa, sien : lesquels tous sont possessifs.
\end{lecture}
\origpage{74}{060}
Ce raisonnement permet d’inclure les formes \textit{cestuy cy} (celui-ci) et \textit{cestuy la} (celui-là), qui, elles-aussi, font aussi référence à la personne à qui l’on parle. Les pronoms relatifs sont sujets à la même analyse :
\begin{lecture}{Les pronoms – Estienne, \textit{Traicté De la grammaire Françoise} (1569)}
Il a été de besoin, pour éviter cette manière de répétition de noms, inventer ces relatifs tant des noms que même des pronoms, comme qui. Car autrement au lieu de dire, Tu es celui qui me plais, il nous eût fallu dire Tu es celui tu me plais et Qui est relatif de Tu.
\end{lecture}
\subsubsection*{4\quad Le participe}
Un participe est un adjectif ainsi nommé parce que les grammairiens anciens le considéraient comme « participant » à la catégorie des noms et des verbes, le participe se combine avec les auxiliaires \textit{être} ou \textit{avoir} pour former des temps composés. Employé comme adjectif, il a la valeur d’un qualificatif et s’accorde en genre et en nombre avec le nom qu’il détermine.

Dans les grammaires du 16\textsuperscript{e} siècle, le participe occupe une position à part entière dans la liste des parties du discours. \textbf{Clément Marot} (1496-1544), poète officiel de François Ier fasciné par la Renaissance italienne, qui en effet a introduit\ednote{原文此句缺主句谓语，« Clément Marot, ... qui en effet a introduit » 应作 « Clément Marot, ... a en effet introduit »，或在句首补 « C’est » 形成强调句。此注仅订正句法，不据此确认有关“引入”过去分词一致规则的历史归属。} en France la règle du participe passé (appelé « participe passif » par La Ramée) utilisée dans la langue italienne. Dans une \textit{Épigramme à ses disciples} (1538), il en a énoncé la règle en vers :
\begin{lecture}{Lecture : \textit{Épigrammes ses disciples}}
\ednote{本引文标题作 « Épigrammes ses disciples »，缺介词 à，并与上一句的单数书目题名不一致；按上一句应作 « Épigramme à ses disciples »。}
\begin{center}
Notre langue a cette façon\\
Que le terme qui va devant\\
Volontiers régit le suivant.\\
Faut dire en termes parfaits :\\
« Dieu en ce monde nous a faits » ;\\
Faut dire en paroles parfaites :\\
« Dieu en ce monde les a faites » ;\\
Et ne faut point dire en effet :\\
« Dieu en ce monde les a fait ».\\
Ni « nous a fait » pareillement,\\
Mais « nous a faits » tout rondement.\\
L’italien, dont la faconde\\
Passe les vulgaires du monde,\\
Son langage a ainsi bâti\\
En disant : Dio ci à fatti.
\end{center}
\end{lecture}
\origpage{75}{061}
Ainsi, on pourrait dire que c’est Marot qui est responsable des migraines occasionnées par l’accord du participe passé.
\subsubsection*{5\quad Prémisses de l’analyse grammaticale}
Au 16\textsuperscript{e} siècle, l’analyse grammaticale est n’est encore que balbutiante.\ednote{原文 « est n’est » 重复，应为 « l’analyse grammaticale n’est encore que balbutiante »。} Meigret propose un système de fonctions qui ne repose que sur deux fonctions : le « surposé » qui correspond à notre sujet, et le « sousposé », qui correspond au complément du verbe. Les notions « d’agent » et de « patient » leur sont déjà associées.
\begin{lecture}{Les fonctions – Meigret, \textit{Le tretté de grammère françoise} (1550)}
J’appelle le nom surposé ou apposé, celui qui gouverne le verbe, et le sousposé ou souposé, celui qui est gouverné ; comme Pierre aime Laurence, là où Pierre est le surposé, et Laurence, le sousposé : ce que ne se doit pas entendre selon l’ordre des paroles, mais selon le sens : car celuy qui gouverne est réputé en verbes actifs, comme agent, et celui qui est gouverné, comme patient.
\end{lecture}
La Ramée, lui aussi, s’est intéressé à l’accord (« la convenance ») du sujet et du verbe, et a proposé deux fonctions : le suppôt (le sujet) et l’appôt (le verbe).
\begin{lecture}{}
Venons maintenant à la convenance du nom et du verbe. La convenance du nom avec le verbe est en nombre et en personne. Le nom précédent devant le verbe est ici appelé suppôt, le verbe appôt.
\end{lecture}
\section{GRAMMAIRES ET GRAMMAIRIENS AU 17\textsuperscript{e} SIÈCLE}
\markright{Grammairiens au 17\textsuperscript{e} siècle}
Le 17\textsuperscript{e} siècle marque un tournant décisif dans l’évolution de la grammaire. D’un côté, on retrouve la grammaire normative que nous connaissons déjà, basée sur le bon usage, et socialement élitiste. Mais au même moment, une grammaire d’un nouveau genre, analytique et descriptive, inspirée par la philosophie et la logique, a émergé. C’est elle qui préfigure la naissance de la linguistique moderne.
\subsection{Les « censeurs de la langue »}
Sous le règne de Louis XIII (1610-1643), le français était essentiellement la langue de la cour, de l’aristocratie et de la bourgeoisie, et était parlé par moins d’un million de personnes sur une population totale de vingt millions. En 1634, le cardinal de Richelieu (1585-1642) a créé l’Académie française, dont la mission, fixée en 1637, était de nettoyer le français des « ordures » qu’il avait contractées dans la bouche du peuple, donner des règles à la langue française, la rendre « pure », « éloquente » et « capable de traiter des arts et des sciences ». Les « ordures » désignaient les mots mal employés, les régionalismes ou les dialectes, les mots étrangers, les termes techniques et les jargons. Dans ce siècle d’organisation autoritaire et centralisée, ce sont les grammairiens qui ont façonné la langue à leur goût et le règne de Louis XIV a produit bon nombre de « censeurs » professionnels.
\origpage{76}{062}
\subsubsection*{1\quad Claude Favre de Vaugelas : « le greffier de l’usage »}
De façon incontestable, le plus célèbre de ces spécialistes de la langue, celui dont l’autorité était la plus grande, est le grammairien \textbf{Claude Favre de Vaugelas} (1585-1659),\ednote{原文卒年 1659 有误，应为 1650。Académie française 的官方传记记载其卒于 1650 年 2 月 26 日，见\ecite{15}。}
élu en 1634 membre de l’Académie française alors qu’il n’avait encore rien publié : ce n’est qu’en 1647 qu’il a publié son principal ouvrage \textit{Remarques sur la langue française, utiles à ceux qui veulent bien parler et bien écrire}. Dans cet ouvrage qui a connu un immense succès, Vaugelas cherche à définir et à codifier le « bon usage » du français en s’inspirant de la langue parlée à la cour du roi :
\begin{lecture}{}
Le mauvais se forme du plus grand nombre de personnes, qui presque en toutes choses n’est pas le meilleur, et le bon au contraire est composé non pas de la pluralité, mais de l’élite des voix, et c’est véritablement celui que l’on nomme le maître des langues. Voici donc comment on définit le bon usage : c’est la façon de parler de la plus saine partie de la Cour.
\end{lecture}
Préoccupés d’« épurer » la langue, Vaugelas et ses disciples ont proscrit les italianismes, les archaïsmes, les provincialismes, les néologismes, les termes techniques et savants, bref tous les mots « bas ».
\begin{bbcomplement}{}
Vaugelas a été immortalisé sous la plume de Molière dans \textit{Les femmes savantes} (1672), ce qui atteste de sa grande notoriété. Aujourd’hui encore, dire d’une personne qu’elle « parle comme Vaugelas », c’est rendre hommage à sa maîtrise parfaite de la langue française.
\minititle{Lecture : \textit{Les femmes savantes} (extrait)}
CHRYSALE.
\par
Vous êtes satisfaite, et la voilà partie.\\
Mais je n’approuve point une telle sortie ;\\
C’est une fille propre aux choses qu’elle fait,\\
Et vous me la chassez pour un maigre sujet.\\
PHILAMINTE.\\
Vous voulez que toujours je l’aie à mon service,\\
Pour mettre incessamment mon oreille au supplice ?\\
Pour rompre toute loi d’usage et de raison,\\
Par un barbare amas de vices d’oraison,\\
De mots estropiés, cousus par intervalles,\\
De proverbes traînés dans les ruisseaux des Halles ?\\
BELISE.\\
Il est vrai que l’on sue à souffrir ses discours :\\
Elle y met \textbf{Vaugelas} en pièces tous les jours ;
\origpage{77}{063}
Et les moindres défauts de ce grossier génie\\
Sont ou le pléonasme, ou la cacophonie.\\
CHRYSALE.\\
Qu’importe qu’elle manque aux lois de \textbf{Vaugelas},\\
Pourvu qu’à la cuisine elle ne manque pas ?\\
J’aime bien mieux, pour moi, qu’en épluchant ses herbes,\\
Elle accommode mal les noms avec les verbes,\\
Et redise cent fois un bas ou méchant mot,\\
Que de brûler ma viande, ou saler trop mon pot.\\
Je vis de bonne soupe, et non de beau langage.\\
\textbf{Vaugelas} n’apprend point à bien faire un potage ;\\
Et \textbf{Malherbe} et \textbf{Balzac}, si savants en beaux mots, »\ednote{本行末有未与开引号配对的闭引号 »；按此处引文的连续结构应删去。保留底本并在此说明，不改变诗句内容。}\\
En cuisine peut-être auraient été des sots.

Une petite précision pour ceux et celles d’entre vous qui seront étonnés de voir Molière évoquer Balzac. Le Balzac de Molière n’est pas Honoré de Balzac mais Jean-Louis Guez de Balzac, poète ayant beaucoup contribué à réformer la langue française et surnommé « le restaurateur de la langue française ». Quant au poète François de Malherbe, il était surnommé le « tyran de syllabes » en raison de son purisme.
\end{bbcomplement}
\subsubsection*{2\quad Dominique Bouhours : le continuateur de Vaugelas}
Prêtre jésuite français, grammairien et historien, \textbf{Dominique Bouhours} (1628-1702) a été le continuateur de Vaugelas. Comme ce dernier, il a été le maître à penser et à écrire de sa génération et a exercé une influence non négligeable sur des auteurs tels que Boileau, La Fontaine, et Racine, ce dernier lui envoyant même ses pièces pour qu’il les corrige.

Dans ses \textit{Entretiens d’Ariste et d’Eugène} (1671), ouvrage qui a connu un grand succès dans toute l’Europe jusqu’à la Révolution, Bouhours écrit :
\begin{lecture}{}
Il n’y a guère de pays dans l’Europe où l’on n’entende le françois et il ne s’en faut rien que je ne vous avoue maintenant que la connaissance des langues étrangères n’est pas beaucoup nécessaire à un François qui voyage. Où ne va-t-on point avec notre langue ?
\end{lecture}
Cette thèse de l’universalité de la langue française et la supériorité de l’esprit français sera reprise dans un de ses ouvrages les plus célèbres, \textit{La Manière de bien penser dans les ouvrages d’esprit} (1687).
\origpage{78}{064}
\begin{bbcomplement}{François de Malherbe : le « tyran des syllabes »}
Proche d’Henri IV et de Louis XIII, poète officiel de 1605 à 1628, \textbf{François de Malherbe} (1555-1628) était, comme l’a écrit Molière, « savant en beaux mots » et a fait de l’épuration de la langue française l’œuvre de sa vie. On peut le considérer comme le premier théoricien de l’art classique et l’un des réformateurs de la langue française. Durant tout le 17e siècle, il a été la référence majeure des théoriciens classiques.

Sa doctrine consistait, selon ses propres mots, à « mettre le holà » (c’est à dire mettre fin) à des changements incessants. Les principes que Malherbe recherchait dans sa doctrine étaient : « [la] correction, [la] clarté, [la] plénitude, [l’] harmonie, [et la] sobriété de la forme ». Chaque mot pour Malherbe ne pouvait être employé que dans un contexte spécifique. C’est le point principal de sa doctrine : rechercher toutes les fautes du mauvais usage dans la langue française.

Partageons ici un des très rares poèmes de Malherbe ayant traversé les siècles : la \textit{Consolation à M. Du Périer} (1598), qui avait perdu sa fille en très bas âge :
\minititle{\textit{Consolation à M. Du Périer} (1598)}
\begin{center}
Ta douleur, Du Périer, sera donc éternelle,\\
Et les tristes discours\\
Que te met en l’esprit l’amitié paternelle\\
L’augmenteront toujours !\\
Le malheur de ta fille au tombeau descendue\\
Par un commun trépas,\\
Est-ce quelque dédale, où ta raison perdue\\
Ne se retrouve pas ?\\
Je sais de quels appas son enfance était pleine,\\
Et n’ai pas entrepris,\\
Injurieux ami, de soulager ta peine Avecque son mépris.\\
Mais elle était du monde où les plus belles choses\\
Ont le pire destin, Et rose elle a vécu ce que vivent les roses,\\
L’espace d’un matin.
\end{center}
\end{bbcomplement}
\subsection{\textit{La grammaire} de Port Royal : le grand tournant}
Fondée en 1204 au sud-ouest de Paris, l’abbaye de Port-Royal des Champs a été pendant longtemps le théâtre d’une intense vie religieuse, intellectuelle et politique. En raison de son isolement, l’abbaye était qualifiée d’« affreux désert » par la \textbf{marquise de Sévigné} (1626-1996).\ednote{原文卒年 1996 为排印错误，应为 1696。Hachette Livre–BnF 重印书信集的书目资料明确标示 1626–1696，见\ecite{16}。}
En raison aussi de cet isolement, les résidents de l’abbaye étaient appelés les « solitaires ». Parmi les plus célèbres « Solitaires » de Port Royal, on compte Blaise Pascal ou Jean Racine, qui y a été scolarisé dès son plus jeune âge.
\origpage{79}{065}
Au 17e siècle, Port Royal était le symbole de la contestation politique et religieuse face à l’absolutisme royal et aux abus de l’Église. Les jésuites, par l’intermédiaire de Louis XIV ont fini par chassé\ednote{原文 « ont fini par chassé » 应为 « ont fini par chasser »：finir par 后接不定式。} les religieux en 1709 avant de faire raser l’abbaye en 1712. L’abbaye de Port-Royal des Champs était un lieu de réunion où de grands écrivains, théologiens, mathématiciens, philosophes ont produit nombre de travaux sur la logique, la grammaire, la théologie, ainsi que des traductions des textes religieux. Le nom de Port Royal reste de nos jours attaché à deux ouvrages célèbres : la \textit{Grammaire} (1660) et la \textit{Logique} (1662).
\subsubsection*{1\quad D’une grammaire « générale et raisonnée » à une grammaire universelle}
En 1660, le prêtre, théologien, philosophe et mathématicien \textbf{Antoine Arnauld} (1612-1694) et le grammairien \textbf{Claude Lancelot} (1615-1695), ont publié une grammaire, dont le titre original est \textit{Grammaire générale et raisonnée contenant les fondements de l’art de parler expliqués d’une manière claire et naturelle}. Fondée non pas sur le bon usage, mais sur la raison (le \textit{logos}), c’est l’antithèse de la grammaire normative de Vaugelas. Directement inspirée des \textit{Règles pour la direction de l’esprit} (1628-1629) de \textbf{Descartes} pour qui la raison est universelle, elle s’intéresse au langage en général, et non à une langue déterminée. En cela, elle devrait vous rappeler la grammaire des modistes présentée un peu plus haut. Croyant en l’existence de mécanismes logiques universels que chaque langue exprime dans son propre système, la \textit{Grammaire} de Port Royal a profondément marqué la naissance de la linguistique moderne.

Descartes, dans son \textit{Discours de la méthode}, affirme que la raison un\ednote{原文 « la raison un » 缺系词，应为 « la raison est un »。} « instrument universel, qui peut servir en toutes sortes de rencontres ». Suite à Port-Royal, l’idée d’universalité de la grammaire a inspiré au linguiste allemand \textbf{Wilhelm von Humboldt} (1767-1835) sa célèbre théorie de la « forme intérieure du langage », niveau profond du langage totalement indépendant des variations nationales et individuelles. Plus tard encore, on pourra trouver une forte analogie entre le principe de l’universalité des structures du langage défendu par Port-Royal et les hypothèses des idées innées (\textit{innate ideas}), des universaux du langage (\textit{language universals}) et de structure profonde (\textit{deep structure}) formulées par Noam Chomsky.
\subsubsection*{2\quad Dualité du signe linguistique}
S’inspirant du dualisme cartésien du corps et de l’âme présenté dans le chapitre précédent, les grammairiens de Port-Royal soulignent la dichotomie (c’est-à-dire la séparation en deux composantes) caractérisant toutes les manifestations linguistiques. Cette dualité est la marque distinctive de tout signe linguistique. Dans \textit{Les mots et les choses}, \textbf{Michel Foucault} (1926-1984) cite explicitement la \textit{Logique} de Port-Royal (1662), qui définit ainsi le signe :
\begin{lecture}{}
« Ainsi le signe enferme deux idées, l’une de la chose qui représente, l’autre de la chose représentée, et sa nature consiste à exciter la seconde par la première. »

Le « son » est la face extérieure et apparente du signe, son incarnation physique ou matérielle. Le « sens », quant à lui, constitue le côté intérieur, mental du signe.
\end{lecture}
\origpage{80}{066}
\subsection{L’héritage de Port Royal}
Dans le prochain chapitre, consacré à la logique et à la philosophie du langage, nous verrons qu’un énoncé est constitué par un système de propositions élémentaires formées d’un sujet et d’un prédicat, comme dans « Socrate (sujet) est mortel (prédicat) ». Cette forme de jugement logique est, selon Port-Royal, attestée universellement dans toutes les langues. On comprend dès lors plus aisément pourquoi Chomsky a trouvé dans la \textit{Grammaire} de Port-Royal l’inspiration de sa \textit{Grammaire Générative}, dont le but est de faire apparaître, au-delà des structures de surface, les structures profondes.
\subsubsection*{1\quad \textit{La linguistique cartésienne} de Noam Chomsky}
Dans \textit{La linguistique cartésienne} (1969), Chomsky reconnaît à plusieurs reprises sa filiation avec la \textit{Grammaire} de Port Royal :
\begin{lecture}{Lecture : \textit{La linguistique cartésienne}, extrait (1)}
Pour user d’une terminologie récente, nous pouvons distinguer « la structure profonde » d’une phrase de sa « structure de surface ». La première est la structure abstraite et sous-jacente qui détermine l’interprétation sémantique ; la seconde est l’organisation superficielle d’unités qui détermine l’interprétation phonétique et qui renvoie à la forme physique de l’énoncé effectif, à sa forme voulue ou perçue.
\end{lecture}
Puis Chomsky évoque de façon non équivoque sa dette envers la \textit{Grammaire} de Port-Royal.
\begin{lecture}{Lecture : \textit{La linguistique cartésienne}, extrait (2)}
Apparemment, la grammaire de Port-Royal est la première à développer de façon à peu près claire la notion de structure syntagmatique. Il est donc intéressant de remarquer qu’elle établit très clairement l’inadéquation de la description syntagmatique pour représenter la structure syntactique, et qu’elle fait allusion à une forme de grammaire transformationnelle qui est, à bien des égards, proche de celle que nous étudions activement aujourd’hui.
\end{lecture}
\subsubsection*{2\quad Michel Foucault : Port-Royal et l’« épistémè »}
Dans \textit{Les mots et les choses}, Michel Foucault développe la notion d’épistémè, concept philosophique que l’on pourrait traduire par « science », mais dans le double sens de « savoir », et de « vertu », vertu qui consiste à « être savant en acte ». Dans cet essai, Foucault a identifié quelques penseurs qui ont été déterminants dans la mise en place de l’épistémè moderne, et \textit{La Logique} de Port-Royal en fait partie. L’année suivante, en 1967, Foucault a publié dans la revue \textit{Langage} un article intitulé « La grammaire générale de Port Royal ». Il y distingue deux grammaires : un « art de parler » et une « discipline contenant les fondements de cet art ».
\origpage{81}{067}
\begin{lecture}{Lecture : \textit{La grammaire générale de Port Royal} (extrait)}
La grammaire elle non plus n’est pas un « art de bien parler », mais tout simplement un « art de parler » [...] De là une conséquence importante : la grammaire ne saurait valoir comme les prescriptions d’un législateur donnant enfin au désordre des paroles leur constitution et leurs lois ; elle ne saurait être non plus comprise comme un recueil des conseils donnés par un correcteur vigilant. Elle est une discipline qui énonce les règles auxquelles il faut bien qu’une langue s’ordonne pour pouvoir exister. Elle a à définir cette correction d’une langue qui n’est ni son idéal, ni son meilleur usage, ni la limite que le bon goût ne saurait franchir, mais la forme et la loi intérieure qui lui permettent tout simplement d’être la langue qu’elle est. Par le fait même, le sens du mot grammaire se dédouble : il y a une grammaire qui est l’ordre imminent\ednote{本引文中 « imminent » 疑为 « immanent » 之误：前者为“即将发生的”，后者为“内在于……的”，与 « à toute parole prononcée » 相接及前文“内在法则”相应。此为拟校，尚未核到对应原刊页的字形，不将它当作已确证的引文订正。} à toute parole prononcée, et une grammaire qui est la description, l’analyse et l’explication, — la théorie — de cet ordre. La grammaire, c’est la loi de ce que je dis ; et c’est aussi la discipline qui permet de connaître cette loi.
\end{lecture}
C’est sur la \textit{Logique} de Port Royal que s’achève ce troisième chapitre et c’est sur la \textit{Logique} d’Aristote que va s’ouvrir le chapitre suivant, tout entier consacré à la science du \textit{logos}.
\clearpage
\section*{Exercices}
\phantomsection
\addcontentsline{toc}{section}{Exercices}
\markright{Exercices — Chapitre 3}
\minititle{I.\quad Quelles différences faites-vous entre les parties du discours en français moderne et en mandarin ? En raison de l’absence d’articles, comment s’effectue la détermination en chinois ?}
\minititle{II.\quad Dans cet extrait de son \textit{Art Poétique} (1674), Nicolas Boileau évoque plusieurs défenseurs de la langue française présentés dans ce chapitre. Quels sont-ils ? Qui est François Villon ? Comment comprenez-vous les expressions suivantes : « Sa muse, en françois parlant grec et latin », « la langue réparée », « l’oreille épurée » ?}
\begin{adjustwidth}{4mm}{4mm}\small\setlength{\parindent}{0pt}
\textbf{Villon} sut le premier, dans ces siècles grossiers,\\
Débrouiller l’art confus de nos vieux romanciers.\\
Marot bientôt après fit fleurir les ballades ;\\
Tourna des triolets, rima des mascarades,\\
A des refrains réglés asservit les rondeaux,\\
Et montra pour rimer des chemins tout nouveaux.\\
\textbf{Ronsard}, qui le suivit, par une autre méthode,
\origpage{82}{068}
Réglant tout, brouilla tout, fit un art à sa mode,\\
Et toutefois longtemps eut un heureux destin.\\
Mais sa muse, \textbf{en françois parlant grec et latin},\\
Vit dans l’âge suivant, par un retour grotesque,\\
Tomber de ses grands mots le faste pédantesque.\\
Ce poète orgueilleux, trébuché de si haut,\\
Rendit plus retenus \textbf{Desportes} et Bertaut.\\
Enfin \textbf{Malherbe} vint, et, le premier en France,\\
Fit sentir dans les vers une juste cadence,\\
D’un mot mis en sa place enseigna le pouvoir,\\
Et réduisit la muse aux règles du devoir.\\
Par ce sage écrivain la \textbf{langue réparée}\\
N’offrit plus rien de rude à l’\textbf{oreille épurée}.\\
Tout reconnut ses lois ; et ce guide fidèle\\
Aux auteurs de ce temps sert encor de modèle.\\
Marchez donc sur ses pas ; aimez sa pureté,\\
Et de son tour heureux imitez la clarté.
\end{adjustwidth}
\minititle{III.\quad À partir de cet extrait des \textit{Remarques sur la langue française}, rapprochez la notion de « mauvais usage » chez Vaugelas de celle de « faute » chez Henri Frei.}
\begin{adjustwidth}{4mm}{4mm}\small\setlength{\parindent}{0pt}
\textbf{Il y a sans doute deux sortes d’usages, un bon et un mauvais.} Le mauvais se forme du plus grand nombre de personnes, qui presque en toutes choses n’est pas le meilleur, et le bon au contraire est composé non pas de la pluralité, mais de l’élite des voix, et c’est véritablement celui qu’il faut suivre pour bien parler, et pour bien écrire en toutes sortes de styles [...]. Voici donc comme on définit le bon Usage. C’est la façon de parler de la plus saine partie de la Cour, conformément à la façon d’écrire de la plus saine partie des Auteurs du temps. [...]
\end{adjustwidth}
\section*{Pour aller plus loin}
\phantomsection
\addcontentsline{toc}{section}{Pour aller plus loin}
\markright{Pour aller plus loin — Chapitre 3}
\begingroup\small\setlength{\parindent}{0pt}\setlength{\parskip}{.55em}
ARNAULD (Antoine) et LANCELOT (Claude). \textit{Grammaire générale et raisonnée}. BiblioBazaar, 2009.

ARNAULD (Antoine) et LANCELOT (Claude). \textit{La logique ou l’art de penser}. Édition critique par Dominique Descotes, Paris, Honoré Champion, 2011.\ednote{此书共同作者应为 Antoine Arnauld 与 Pierre Nicole，而非 Claude Lancelot。与本条相同的 Dominique Descotes 2011 年校订版在 CiNii Books 的责任者记录为 « Antoine Arnauld, Pierre Nicole »，见\ecite{17}。Lancelot 是前一条《Grammaire générale et raisonnée》的共同作者。}

CHEVALIER (Jean-Frédéric) et CAPELLA (Martianus). \textit{Les noces de Philologie et de Mercure}. Livre I, Paris, \origcont{83}{069}Les Belles Lettres, 2014.

CHOMSKY (Noam). \textit{La Linguistique cartésienne suivi de La Nature formelle du langage}. Paris, Seuil, 1969.

DONATUS (Aelius). \textit{Ars Grammatica}. Turnhout, Brepols, 1982.

LESAULNIER (Jean) et MCKENNA (Antony). \textit{Dictionnaire de Port-Royal}. Honoré Champion, 2004.

Pāṇini. \textit{The Aṣṭādhyāyī of Pāṇini with translation and explanatory notes}. Sahitya Akademi, 1992.

ROSIER (Irène). \textit{La grammaire spéculative des modistes}. Presses universitaires de Lille, 1980.

VAUGELAS (René L. de).\ednote{此条把 Vaugelas 的名字写为 « René L. »，与正文及章首提纲的 « Claude Favre de Vaugelas » 不一致；作者名应订正为 « VAUGELAS (Claude Favre de) »。此注只订正作者姓名，不据此认定所列 1975 年版本的全部出版信息无误。} \textit{Remarques sur la langue française}. Larousse, 1975.

WADDINGTON (Charles-Pendrell). \textit{Pierre (Pierre de la Ramée) : sa vie, ses écrits et ses opinions}. Hachette, 2013.
\endgroup
\section*{Glossaire}
\phantomsection
\addcontentsline{toc}{section}{Glossaire}
\markright{Glossaire — Chapitre 3}
\gloss{descriptif 描述的，描写的}{在语法中，描写语法描写语言在实际中怎么说怎么写，而不是说明或规定该怎么说或怎么写。}
\gloss{grammaire 语法}{对语言结构及语言单位如单词、短语在该语言中组成句子方式的描写，通常包括这些句子在整个语言体系中的意义及功能。}
\gloss{grammaire universelle 普遍语法}{这种理论认为，不管操什么语言，每个成人的语法都可用这种理论来描述。}
\gloss{morphologie 形态学，构词学，词法学}{研究词素、词素的不同形式（即词素变体）及其在构词时的组合方式的一门学科。}
\gloss{phonétique 语音的，语音学}{研究语音的科学，包括三大研究范围：发音语音学、声学语音学、听觉语音学。}
\gloss{phonologie 音系学，音位学}{指研究或描述一种语言里有区别性的语音单位（音位）及其相互关系的学科。也可表示确定一种语言中音位的方法。}
\origpage{84}{070}
\gloss{prescriptif 规定的}{在语法中，规定语法用来阐述那些被认为是最好的或者最正确用法规则的一种语法。}
\gloss{sémantique 语义学，语义学的}{关于意义的研究。研究语言的意义有多种不同方法，如哲学家研究语言表达法之间的关系；语言学家研究语言的语义结构，并区分不同的意义种类。}
\gloss{signe 符号}{语言学中，指代表事物的单词或其他词语。有些语言学家和哲学家给符号化过程加入了第三项，即符号所代表的事物的抽象概念。}
